% gnn_transformer_diffused_attention.tex
% Compile with: pdflatex gnn_transformer_diffused_attention.tex
\documentclass[tikz,border=6pt]{standalone}
\usepackage{tikz}
\usetikzlibrary{positioning,arrows.meta,shapes.geometric,shadows,fit,calc,decorations.pathreplacing}
\usepackage{amsmath}
\usepackage{siunitx}
\usepackage{varwidth}

% --- Styling parameters (tweak these) ---
\def\scale{1.0}
\definecolor{colProtein}{RGB}{90,155,210}
\definecolor{colEdge}{RGB}{80,120,160}
\definecolor{colGNN}{RGB}{106,168,79}
\definecolor{colTrans}{RGB}{66,133,244}
\definecolor{colDiff}{RGB}{219,68,55}
\definecolor{colProj}{RGB}{171,71,188}
\definecolor{boxbg}{RGB}{250,250,250}
\tikzset{
  arr/.style={-{Latex[length=3mm,width=2mm]}, line width=0.9pt},
  arrdashed/.style={-{Latex[length=3mm,width=2mm]}, dashed, line width=0.8pt},
  module/.style={draw=black, fill=boxbg, rounded corners=2pt, inner sep=6pt, minimum width=36mm},
  smallmod/.style={draw=black, fill=boxbg, rounded corners=2pt, inner sep=3pt, font=\footnotesize, minimum width=28mm},
  circ/.style={circle, draw=black, fill=white, minimum size=7mm, inner sep=0pt, font=\small},
  token/.style={rectangle, draw=black, rounded corners=1pt, minimum width=8mm, minimum height=6mm, inner sep=1pt},
  lbl/.style={font=\footnotesize},
  eqn/.style={font=\footnotesize\itshape}
}

\begin{document}
\begin{tikzpicture}[scale=\scale, every node/.style={transform shape}]

% -----------------------------
% Title (for draft; remove for final)
% -----------------------------
\node[anchor=west,font=\bfseries] at (-9.8,6.8) {\large GNN \(\rightarrow\) Transformer with Diffusively Gated Attention};

% =============================
% Left: 3D Protein / geometric graph
% =============================
% stylized 3D protein cloud built from spheres
\begin{scope}[shift={(-8.6,2)}]
  % protein backbone spheres (3D impression by vertical offsets)
  \foreach \i/\x/\y/\z in {
    1/-0.9/0.2/0.6, 2/-0.4/0.7/0.9, 3/0.4/0.9/0.7, 4/1.0/0.6/0.3,
    5/0.6/0.0/0.2, 6/0.0/-0.5/0.4, 7/-0.7/-0.7/0.6, 8/-1.2/-0.1/0.5, 9/-0.2/0.3/1.1
  }{
    \shade[ball color=colProtein!60!white] (\x,\y) circle (0.42cm);
    \node[circ, minimum size=6mm, fill=colProtein!10] at (\x,\y) {};
  }
  % edges for structure
  \draw[line width=0.8pt, color=colEdge] (-0.9,0.2) -- (-0.4,0.7) -- (0.4,0.9) -- (1.0,0.6) -- (0.6,0.0) -- (0.0,-0.5) -- (-0.7,-0.7) -- (-1.2,-0.1) -- (-0.9,0.2);
  \draw[line width=0.7pt, color=colEdge] (-0.4,0.7) -- (-0.2,0.3);
  % label
  \node[lbl, anchor=north] at (0.2,-1.05) {3D protein structure (PDB)};
  % arrow to conversion
  \draw[arr] (1.3,0.6) -- (2.3,0.6) node[midway, above, font=\scriptsize] {construct geometric graph};
\end{scope}

% =============================
% Graph / geometric graph representation (nodes as 3D coords)
% =============================
\node[module, anchor=west] (graphbox) at (-3.4,2.5) {\begin{minipage}{44mm}
  \centering \textbf{Geometric Graph}\\
  $G=(V,E)$ with node coords $\{p_v\in\mathbb{R}^3\}$ and node features $x_v$ \\
  (distance / angle / torsion features, residue identity)
\end{minipage}};

% draw little graph inside/near the module
\begin{scope}[shift={(-3.7,1.5)}]
  \foreach \i/\x/\y in {1/0/0.5,2/0.7/1.0,3/1.5/0.2,4/0.6/-0.5,5/-0.4/-0.3}{
    \node[circ, fill=colProtein!6] (gv\i) at (\x,\y) {$v_{\i}$};
  }
  \draw[line width=0.7pt] (gv1) -- (gv2) -- (gv3) -- (gv5) -- (gv4) -- (gv1);
  \draw[line width=0.6pt] (gv2) -- (gv4);
\end{scope}

% arrow to GNN
\draw[arr] (-1.3,2.5) -- (-0.1,2.5) node[midway, above,font=\scriptsize] {message passing};

% =============================
% GNN encoder with geometric attention
% =============================
\node[module, anchor=west, fill=colGNN!8] (gnn) at (0.6,2.0) {\begin{minipage}{54mm}
  \centering \textbf{GNN Encoder (geometric attention)}\\[2pt]
  Per-layer message passing with geometric attention:\\
  \(\displaystyle m_{v}^{(t)}=\sum_{u\in\mathcal{N}(v)} \alpha_{uv}^{(t)}\; W_m h_u^{(t)},\quad
    \alpha_{uv}^{(t)}=\mathrm{softmax}_u\big(f(h_u,h_v,\phi(p_u,p_v))\big)\)\\[3pt]
  Output: node embeddings \(H=[h_v]_{v\in V}\)
\end{minipage}};

% small stack indicator
\node[smallmod, below=of gnn, yshift=8pt] (gnnstack) {Stack of $L_{GNN}$ geometric attention layers};

% arrow: embeddings to projection
\draw[arr] (2.0,1.3) -- (3.2,1.3) node[midway, above, font=\scriptsize] {$H\in\mathbb{R}^{n\times d}$};

% optional cross-link (GNN features -> Transformer K/V)
\draw[arrdashed] (1.9,1.85) .. controls (2.6,1.9) and (3.2,2.7) .. (4.0,3.7) ;
\node[lbl] at (3.25,2.55) {use as K/V (optional)};

% =============================
% Graph-to-sequence projection (learnable tokens / pooling)
% =============================
\node[smallmod, anchor=west, fill=colProj!10] (proj) at (3.6,1.0) {\begin{minipage}{40mm}\centering
  Graph\(\to\)Sequence Projection \\ (learned readout tokens or pooling) \\
  \(S = \mathrm{Proj}(H) = [\tau_1,\ldots,\tau_m]\)
\end{minipage}};

% tokens
\foreach \i in {1,2,3,4}{
  \node[token, anchor=west] (tok\i) at ($(proj.east)+(6mm + 10mm*(\i-1),0.0)$) {$\tau_{\i}$};
}
\draw[arr] (proj.east) -- (tok1.west);

% labels for sequence (mutation info fed later)
\node[lbl, anchor=west] at (6.9,0.2) {sequence tokens};

% =============================
% Sequence Transformer with mutation-aware attention
% =============================
\node[module, anchor=west, fill=colTrans!7] (trans) at (8.6,1.0) {\begin{minipage}{72mm}
  \centering \textbf{Sequence Transformer (mutation-aware)}\\[2pt]
  Standard multihead self-attention + mutation conditioning:\\[2pt]
  \(\displaystyle A=\mathrm{softmax}\!\big(\tfrac{QK^\top}{\sqrt{d_k}}\big),\quad 
    \widetilde A = \Psi_{\text{mut}}(A,M)\)\\
  where \(M\) indicates mutation site / allele info (learned embedding)
\end{minipage}};

% arrow from tokens to transformer
\draw[arr] ($(tok2.east)+(0,0)$) -- ($(trans.west)+(-9mm,0)$);

% small note for mutation-aware attention
\node[lbl] at (11.6,2.3) {mutation embeddings $M$ join Q/K/V / bias attention};

% =============================
% Diffusively gated attention fusion (visual center)
% =============================
% central block describing the fusion formula
\node[module, anchor=west, fill=colDiff!6] (fusion) at (8.6,-2.2) {\begin{minipage}{74mm}
  \centering \textbf{Diffusively Gated Attention Fusion}\\[2pt]
  Combine Transformer attention and structural prior via a differentiable gate: \\[4pt]
  \(\displaystyle A_{\text{struct}}=\Phi(G)\quad(\text{e.g. heat kernel, PPR, }e^{-\tau L})\)\\[4pt]
  \(\displaystyle g = \sigma\big(w_g^\top[\mathrm{pool}(H),\,\mathrm{pool}(Y)] + b_g\big)\in[0,1]\)\\[4pt]
  \(\displaystyle A_{\text{fused}} = g\odot A_{\text{struct}} + (1-g)\odot \widetilde A\)\\[4pt]
  Attention output: \(O = A_{\text{fused}} V\)
\end{minipage}};

% arrows: structural prior from graph to fusion
\draw[arrdashed] ($(graphbox.east)+(0,0.4)$) .. controls ($(graphbox.east)+(1.6,0.8)$) and ($(fusion.west)+( -18mm, 0.8)$) .. ($(fusion.west)+( -4mm, 0.7)$)
 node[pos=0.62, above, font=\scriptsize] {compute $\Phi(G)$};

% arrows: Transformer attention to fusion
\draw[arr] ($(trans.south)+(0,-0.6)$) -- ($(fusion.north)+(0,0.5)$) node[midway,right,font=\scriptsize] {$\widetilde A$ / $Y$};

% Labeled gate
\node[smallmod, fill=colDiff!14, anchor=north] (gate) at ($(fusion.north)+(2.8,-0.9)$) {\begin{minipage}{28mm}\centering
  Gate $g$\\ \(\sigma(\cdot)\) \end{minipage}};
\draw[arr] (gate.west) -- ($(fusion.north)+( -6.8,-0.6)$);

% =============================
% Downstream outputs: Prediction head + reasoning modules
% =============================
\node[module, anchor=west, fill=boxbg] (pred) at (13.2,-2.2) {\begin{minipage}{56mm}
  \centering \textbf{Prediction Head(s)}\\
  - Stability / activity regression (MLP)\\
  - Function classification (softmax)\\
  - Per-residue effect (token-wise head)
\end{minipage}};

\draw[arr] ($(fusion.east)+(0,0)$) -- ($(pred.west)$) node[midway, above,font=\scriptsize] {fused token embeddings / pooled vector};

% reasoning module (example: counterfactual / attribution / pathway)
\node[smallmod, anchor=west, fill=boxbg] (reason) at (13.2,0.6) {\begin{minipage}{56mm}\centering
  \textbf{Reasoning / Explainability}\\
  - Attention attribution (A_fused)\\
  - Counterfactual mutation simulation\\
  - Local structural reweighting
\end{minipage}};
\draw[arr] ($(trans.east)+(5mm,0.6)$) -- ($(reason.west)+(0,0.0)$) node[midway, above,font=\scriptsize] {token embeddings / attention maps};

% final outputs arrows to a small result panel
\node[rectangle, draw=black, rounded corners=2pt, fill=white, minimum width=28mm, minimum height=18mm] (result) at (16.0,-4.5) {\begin{varwidth}{6.2cm}
  \vspace{1pt}\centering \textbf{Final outputs}\\
  Predicted ΔActivity, ΔStability, Important residues, Explanations
  \vspace{1pt}
\end{varwidth}};
\draw[arr] (pred.east) -- (result.west);
\draw[arr] (reason.south) .. controls +(0,-0.8) and +(-0.6,0.4) .. (result.north);

% =============================
% Legend & small equations
% =============================
% Put a compact legend on bottom-left
\node[draw=none, anchor=west] at (-9.8,-3.6) {\begin{minipage}{7.8cm}
  \vspace{0.8mm}\textbf{Legend / Notes:}\\[-2pt]
  {\scriptsize
  \begin{tabular}{@{}l@{}}
   \(\Phi(G)\): structural diffusion / prior matrix (e.g.\ PPR, heat kernel \(e^{-\tau L}\)).\\
   \(\widetilde A\): transformer attention (mutation-aware).\\
   Gate \(g\) can be scalar, vector, or token-wise; learnable end-to-end.\\
   Geometric attention \(f(h_u,h_v,\phi(p_u,p_v))\) encodes distances, angles, torsions.\\
  \end{tabular}
  }
\end{minipage}};

% small decorative brace and label for overall flow
\draw[decorate, decoration={brace, amplitude=6pt}, thick] (-10.6,3.6) -- (5.8,3.6) node[midway, above=8pt,font=\small] {Structural path (left) \(\rightarrow\) representation \(\rightarrow\) Transformer fusion (center)};
\draw[decorate, decoration={brace, amplitude=6pt}, thick] (5.5,-0.6) -- (16.8,-0.6) node[midway, above=8pt,font=\small] {Transformer path (top-right) \(\rightarrow\) diffused fusion \(\rightarrow\) outputs};

% =============================
% Fine-tuning hints (not printed in final)
% =============================
% (commented — left here for authors)
% - To show token-to-node alignment lines, draw dashed lines from tokens to graph nodes.
% - To display multi-block Transformer stack, replicate 'trans' vertically or add "xN" label.
% - Replace the schematic 3D protein with an exported PDB cartoon image if desired in the paper.
% - Use font and color adjustments to match journal style.

\end{tikzpicture}
\end{document}
